\section{Critical Assessment of Alvi's Claimed Contributions}
\label{sec:alvi-critique}

Alvi's dissertation (Chapter 5, Conclusions \& Future Work)
\autocite{alvi2018pgm} closes with eight claimed contributions. Each is
assessed below against what this project was actually able to replicate
and observe first-hand, not merely against the original text's own
framing.

\begin{enumerate}
  \item \textbf{Replacing EGARCH-M-derived views with Bayesian-model-
    derived views for Black-Litterman.} \emph{Achieved, and now
    demonstrated end-to-end in Section~\ref{sec:results} of this
    report.} Alvi's own dissertation proposes this substitution as
    future work rather than implementing it; our
    \texttt{BlackLittermanAllocator} demonstration is, to our knowledge,
    a working instantiation of exactly this idea. It matters because it
    replaces a purely statistical volatility-forecasting input (EGARCH-M)
    with a structurally-interpretable, auditable one (a Bayesian
    network's posterior), directly addressing the "black box" objection
    a risk committee would reasonably raise against a pure time-series
    model.
  \item \textbf{Using HMM-discretized time series as inputs to belief
    networks is novel.} \emph{Partially achieved and independently
    confirmed useful, but "novel" is an overclaim.} We independently
    re-derived the same two-stage design (regime-detect, then learn
    structure over the discretized regimes) without reference to Alvi's
    implementation details, and found it necessary for the same reason
    Alvi likely did: a belief network cannot consume continuous,
    non-stationary price levels directly. However, combining HMM
    state-decoding with downstream structure learning predates this
    2018 dissertation in the broader regime-switching literature; the
    genuine contribution is the specific \emph{application} to crude oil
    risk management, not the general technique.
  \item \textbf{A working, directly-deployable trading mechanism capable
    of independent decision-making.} \emph{Not achieved, by our own
    replication.} Section~\ref{sec:results}'s honest finding --
    69.7\% one-month-ahead regime accuracy, below the 87.9\%
    naive-persistence baseline -- directly contradicts "deployable"; a
    mechanism that underperforms the trivial baseline cannot be handed
    an independent capital-allocation mandate without a human overseeing
    it. This matters because it is exactly the kind of overclaim a risk
    manager must be trained to catch before signing off on a model.
  \item \textbf{Structure-learning-based, so near-autonomous and
    requiring no prior expert knowledge beyond dataset selection.}
    \emph{Achieved, with one caveat.} Our Hill-Climbing pipeline requires
    no manually-specified edges, confirming this claim; the one
    qualification is that \emph{dataset selection itself} is not a
    trivial, expert-knowledge-free step -- deciding to discretize
    financial series via regime-detection but macro/geopolitical series
    via simple terciles (Section 4 of the GWP2 report) is itself a
    domain judgement call that materially shapes what the algorithm can
    possibly discover.
  \item \textit{(Duplicate of contribution 2 in the source text,
    verbatim -- not independently assessed.)}
  \item \textbf{Python PGM libraries' abstraction level aided theoretical
    understanding of the design process.} \emph{Achieved, and we concur
    directly.} Implementing the four HMM algorithms from scratch (GWP2)
    while using \texttt{pgmpy}'s higher-level Hill-Climbing/BIC/CPD
    abstractions for the belief-network side gave us the same
    pedagogical contrast Alvi describes: writing Baum-Welch by hand
    forces an understanding of \emph{why} EM must be used (Section
    "Hidden Markov Models: A Latent-Variable-Model Definition" of the
    GWP2 report), while \texttt{pgmpy}'s \texttt{HillClimbSearch} lets
    attention shift to interpreting the \emph{result} (edge economic
    plausibility, Markov blanket contents) rather than the search
    mechanics.
  \item \textbf{Provides forecasts via a systematic, event-driven
    global-macro strategy with higher-return potential than
    hedge-fund/fixed-income strategies.} \emph{Not supported by our
    replication.} This is the single most over-reaching claim in the
    list: no backtested, cost-adjusted return series is presented
    anywhere in the original dissertation to substantiate a
    hedge-fund-beating return claim, and our own forecast-accuracy
    result (Section~\ref{sec:results}) actively argues against treating
    the model's regime calls as a standalone systematic strategy without
    further validation.
  \item \textbf{Helps energy policy makers understand oil-market
    structure to draft better policy.} \emph{Plausibly achieved, and
    the strongest-surviving claim in our view.} Regardless of the
    model's forecast accuracy, the \emph{structure} it recovers --
    which variables belong in \texttt{wti\_regime}'s Markov blanket, and
    which do not -- is itself informative to a policy audience: our
    replication across three different samples consistently finds that
    an industrial-production indicator and the Brent benchmark, not
    pure financial-market variables, are the proximate drivers of the
    WTI regime. A policymaker monitoring physical-market indicators
    rather than financial-market sentiment to anticipate oil-price
    regime shifts is a defensible, model-structure-derived
    recommendation that does not depend on the forecast-accuracy
    shortfall documented above.
\end{enumerate}

In summary, this project's replication supports Alvi's \emph{structural}
claims (contributions 1, 2, 4, 6, 8) considerably more than his
\emph{performance} claims (contributions 3 and 7): the belief-network
machinery reliably recovers an interpretable, economically plausible,
and reasonably stable causal structure without expert input, but the
resulting model's one-month-ahead forecast accuracy does not, on this
replication, clear the bar of a naive persistence baseline -- a result
future work (and genuine commercial deployment) would need to address
directly rather than assume away.
